%! AP Statistics — Unit 1, Lesson 1.10 — Slides (Beamer)
\documentclass[aspectratio=169, 11pt]{beamer}
\usepackage{apstats-beamer}

%=============================================================================
\begin{document}

% ─────────────────────────────────────────────────────────────────────────────
% SLIDE 1 — LESSON TITLE
% ─────────────────────────────────────────────────────────────────────────────
{
\setbeamercolor{background canvas}{bg=navy}
\begin{frame}[plain]
  \begin{tikzpicture}[remember picture, overlay]
    \fill[navylight] (current page.north west)
      rectangle ([yshift=-1.35cm]current page.north east);
  \end{tikzpicture}
  \vspace{2.8cm}
  \hspace{0.55cm}%
  \begin{minipage}{0.9\textwidth}
    {\color{goldacc}\bfseries\small LESSON 1.10}\\[0.5cm]
    {\color{white}\bfseries\LARGE The Normal Distribution\\[0.25cm]
      and $z$-Scores}\\[1.4cm]
    {\color{skymid}\small AP Statistics~$\cdot$~Unit 1: Exploring One-Variable Data}
  \end{minipage}
\end{frame}
}

% ─────────────────────────────────────────────────────────────────────────────
% SLIDE 2 — WARM-UP
% ─────────────────────────────────────────────────────────────────────────────
{
\setbeamercolor{background canvas}{bg=sky}
\begin{frame}[t, plain]
  \navyheader{Warm-Up}
  \vspace{1em}%
  \hspace{0.55cm}%
  \begin{minipage}{0.9\textwidth}
    \small
    \textit{Use the Normal curve on your warm-up sheet showing adult male heights
    $N(70, 3)$ inches.}

    \vspace{0.5em}
    \begin{enumerate}
      \setlength{\itemsep}{0.6em}\setlength{\topsep}{0em}
      \item What is $\mu$? What is $\sigma$? Write the distribution in $N(\mu,\sigma)$ notation.
      \item The shaded region shows 67 to 73 inches. How many SDs above/below the mean?
      \item Estimate the percentage of men in the shaded region.
      \item A man is 76 inches tall. Is he above or below average? How many SDs away?
    \end{enumerate}
    \vspace{0.4cm}
    {\color{navy}\itshape\small
      Today we build a mathematical model that lets us calculate exact
      probabilities for bell-shaped distributions.}
  \end{minipage}
\end{frame}
}

% ─────────────────────────────────────────────────────────────────────────────
% SLIDE 3 — PRIMARY OBJECTIVE
% ─────────────────────────────────────────────────────────────────────────────
{
\setbeamercolor{background canvas}{bg=sky}
\begin{frame}[t, plain]
  \begin{tikzpicture}[remember picture, overlay]
    \fill[navy] (current page.north west)
      rectangle ([yshift=-0.25cm]current page.north east);
  \end{tikzpicture}
  \vspace{0.45cm}\par
  \hspace{0.55cm}%
  \begin{minipage}{0.9\textwidth}

    \sectionlabel{Primary Objective}

    \normalsize
    Students will describe the properties of the Normal distribution,
    apply the 68--95--99.7 Rule, calculate and interpret $z$-scores,
    use the Normal table to find probabilities and percentiles, and
    assess whether a distribution is approximately Normal.

    \vspace{0.5cm}
    \textbf{Learning Targets:}
    \begin{enumerate}\small
      \setlength{\itemsep}{0.4em}
      \item Describe the Normal distribution; identify $\mu$ and $\sigma$ as parameters.
      \item Apply the Empirical Rule: 68--95--99.7\% within 1, 2, 3 SDs of mean.
      \item Compute $z$-scores; use Table A for $P(Z<z)$, $P(Z>z)$, between.
      \item Find percentiles with the inverse Normal; assess Normality graphically.
    \end{enumerate}
  \end{minipage}
\end{frame}
}

% ─────────────────────────────────────────────────────────────────────────────
% SLIDE 4 — HOOK: THE POWER OF A MODEL
% ─────────────────────────────────────────────────────────────────────────────
{
\setbeamercolor{background canvas}{bg=goldbg}
\begin{frame}[t, plain]
  \navyheader{Hook: The Power of a Mathematical Model}
  \vspace{0.8em}%
  \hspace{0.55cm}%
  \begin{minipage}{0.9\textwidth}
    \small
    Adult male heights $\sim N(70, 3)$ inches. \textbf{Without collecting new data}, we can answer:

    \vspace{0.5em}
    \begin{center}
    \begin{tikzpicture}[scale=0.9]
      \fill[sky, opacity=0.85]
        plot[domain=-1:1, samples=40] (\x, {2.5*exp(-\x*\x/2)})
        -- (1, 0) -- (-1, 0) -- cycle;
      \draw[thick, navy]
        plot[domain=-3.5:3.5, samples=80] (\x, {2.5*exp(-\x*\x/2)});
      \draw[thick] (-3.8, 0) -- (3.8, 0);
      \foreach \z/\h in {-3/61, -2/64, -1/67, 0/70, 1/73, 2/76, 3/79} {
        \draw (\z, -0.10) -- (\z, 0.10);
        \node[below, font=\small] at (\z, -0.13) {\h};
      }
      \node[below, font=\small\itshape] at (0, -0.52) {Height (inches)};
      \node[above, font=\small\bfseries\color{navy}] at (0, 2.52) {$\mu=70$};
    \end{tikzpicture}
    \end{center}

    \vspace{0.3em}
    \begin{itemize}
      \setlength{\itemsep}{0.4em}
      \item What \% of men are between 67 and 73 in.\ (shaded)? \quad $\Rightarrow$ \textbf{68\%}
      \item What \% are taller than 76 in.? \quad $\Rightarrow$ \textbf{2.5\%}
      \item What height is the 90th percentile? \quad $\Rightarrow$ \textbf{use $z$-score}
    \end{itemize}
  \end{minipage}
\end{frame}
}

% ─────────────────────────────────────────────────────────────────────────────
% SLIDE 5 — NORMAL DISTRIBUTION: SHAPE AND PARAMETERS
% ─────────────────────────────────────────────────────────────────────────────
{
\setbeamercolor{background canvas}{bg=white}
\begin{frame}[t, plain]
  \begin{tikzpicture}[remember picture, overlay]
    \fill[navy] (current page.north west)
      rectangle ([yshift=-0.25cm]current page.north east);
  \end{tikzpicture}
  \vspace{0.45cm}\par
  \hspace{0.55cm}%
  \begin{minipage}{0.9\textwidth}

    \sectionlabel{The Normal Distribution}

    \begin{columns}[T]
      \begin{column}{0.52\textwidth}
        \small
        \textbf{Properties:}
        \begin{itemize}
          \setlength{\itemsep}{0.35em}
          \item Symmetric, unimodal, bell-shaped
          \item Mean $=$ Median $=$ Mode (at the peak)
          \item Two parameters: $\mu$ (mean) and $\sigma$ (SD)
          \item Notation: $N(\mu, \sigma)$
          \item Tails approach zero as $x \to \pm\infty$
        \end{itemize}

        \vspace{0.5em}
        \textbf{Effect of parameters:}
        \begin{itemize}
          \setlength{\itemsep}{0.3em}
          \item Changing $\mu$: \textbf{shifts} curve left or right
          \item Larger $\sigma$: curve is \textbf{flatter and wider}
          \item Smaller $\sigma$: curve is \textbf{taller and narrower}
        \end{itemize}
      \end{column}
      \begin{column}{0.44\textwidth}
        \begin{center}
        \begin{tikzpicture}[scale=0.85]
          \draw[thick, navy]
            plot[domain=-3.5:3.5, samples=80] (\x, {2.5*exp(-\x*\x/2)});
          \draw[thick, navy!40]
            plot[domain=-3.5:3.5, samples=80] (\x, {1.5*exp(-\x*\x/8)});
          \draw[thick] (-3.8, 0) -- (3.8, 0);
          \foreach \z in {-3,-2,-1,0,1,2,3} {
            \draw (\z, -0.10) -- (\z, 0.10);
            \node[below, font=\tiny] at (\z, -0.13) {$\z$};
          }
          \node[above, font=\tiny, color=navy] at (0, 2.52) {small $\sigma$};
          \node[right, font=\tiny, color=navy!40] at (2.0, 0.70) {large $\sigma$};
          \node[below, font=\tiny\itshape] at (0, -0.52) {$z$-score};
        \end{tikzpicture}
        \end{center}
      \end{column}
    \end{columns}
  \end{minipage}
\end{frame}
}

% ─────────────────────────────────────────────────────────────────────────────
% SLIDE 6 — THE EMPIRICAL RULE
% ─────────────────────────────────────────────────────────────────────────────
{
\setbeamercolor{background canvas}{bg=white}
\begin{frame}[t, plain]
  \begin{tikzpicture}[remember picture, overlay]
    \fill[navy] (current page.north west)
      rectangle ([yshift=-0.25cm]current page.north east);
  \end{tikzpicture}
  \vspace{0.45cm}\par
  \hspace{0.55cm}%
  \begin{minipage}{0.9\textwidth}

    \sectionlabel{The 68--95--99.7 (Empirical) Rule}

    \begin{columns}[T]
      \begin{column}{0.46\textwidth}
        \small
        In any Normal distribution $N(\mu, \sigma)$:
        \begin{itemize}
          \setlength{\itemsep}{0.4em}
          \item $\approx \mathbf{68\%}$ within $1\sigma$ of $\mu$
          \item $\approx \mathbf{95\%}$ within $2\sigma$ of $\mu$
          \item $\approx \mathbf{99.7\%}$ within $3\sigma$ of $\mu$
        \end{itemize}

        \vspace{0.5em}
        \begin{block}{Example: heights $N(70, 3)$}
          \small
          \begin{itemize}
            \setlength{\itemsep}{0.3em}
            \item 68\%: between 67 and 73 in.
            \item 95\%: between 64 and 76 in.
            \item Above 76: $(100\%-95\%)/2 = 2.5\%$
          \end{itemize}
        \end{block}
      \end{column}
      \begin{column}{0.50\textwidth}
        \begin{center}
        \begin{tikzpicture}[scale=0.88]
          \fill[navy!10]
            plot[domain=-3:3, samples=60] (\x, {2.5*exp(-\x*\x/2)})
            -- (3,0) -- (-3,0) -- cycle;
          \fill[navy!20]
            plot[domain=-2:2, samples=60] (\x, {2.5*exp(-\x*\x/2)})
            -- (2,0) -- (-2,0) -- cycle;
          \fill[sky]
            plot[domain=-1:1, samples=40] (\x, {2.5*exp(-\x*\x/2)})
            -- (1,0) -- (-1,0) -- cycle;
          \draw[thick, navy]
            plot[domain=-3.5:3.5, samples=80] (\x, {2.5*exp(-\x*\x/2)});
          \foreach \z in {-3,-2,-1,1,2,3} {
            \draw[dashed, charcoal, thin] (\z, 0) -- (\z, {2.5*exp(-\z*\z/2)});
          }
          \draw[thick] (-3.8, 0) -- (3.8, 0);
          \foreach \z in {-3,-2,-1,0,1,2,3} {
            \draw (\z, -0.10) -- (\z, 0.10);
            \node[below, font=\tiny] at (\z, -0.13) {$\z$};
          }
          \node[font=\small\bfseries] at (0, 1.0) {68\%};
          \node[font=\small\bfseries] at (1.6, 0.42) {95\%};
          \node[font=\small\bfseries] at (2.6, 0.14) {99.7\%};
          \node[below, font=\tiny\itshape] at (0, -0.52) {$z$-score};
        \end{tikzpicture}
        \end{center}
      \end{column}
    \end{columns}
  \end{minipage}
\end{frame}
}

% ─────────────────────────────────────────────────────────────────────────────
% SLIDE 7 — z-SCORES
% ─────────────────────────────────────────────────────────────────────────────
{
\setbeamercolor{background canvas}{bg=white}
\begin{frame}[t, plain]
  \begin{tikzpicture}[remember picture, overlay]
    \fill[navy] (current page.north west)
      rectangle ([yshift=-0.25cm]current page.north east);
  \end{tikzpicture}
  \vspace{0.45cm}\par
  \hspace{0.55cm}%
  \begin{minipage}{0.9\textwidth}

    \sectionlabel{$z$-Scores: Standardizing a Value}

    \begin{columns}[T]
      \begin{column}{0.50\textwidth}
        \small
        \[
          z = \dfrac{x - \mu}{\sigma}
        \]
        $z$ = number of standard deviations $x$ is above (positive)
        or below (negative) the mean.

        \vspace{0.5em}
        \textbf{Example: SAT scores $N(1060, 195)$}
        \begin{itemize}
          \setlength{\itemsep}{0.4em}
          \item $x = 1255$: $z = \frac{1255-1060}{195} = +1.0$
          \item $x = 900$: $z = \frac{900-1060}{195} \approx -0.82$
          \item $x = 1060$: $z = \frac{1060-1060}{195} = 0$
        \end{itemize}

        \vspace{0.5em}
        Standardizing converts $N(\mu, \sigma)$ to the
        \textbf{Standard Normal} $N(0, 1)$.
      \end{column}
      \begin{column}{0.46\textwidth}
        \begin{block}{Interpretation template}
          \small
          ``A score of 1255 has a $z$-score of $+1.0$,
          meaning it is \textbf{1.0 standard deviation above}
          the mean SAT score of 1060.''
        \end{block}

        \vspace{0.5em}
        \begin{block}{Common errors}
          \small
          \begin{itemize}
            \setlength{\itemsep}{0.3em}
            \item Looking up raw $x$ in Table A
            \item Forgetting the sign of $z$
            \item Not interpreting in context
          \end{itemize}
        \end{block}
      \end{column}
    \end{columns}
  \end{minipage}
\end{frame}
}

% ─────────────────────────────────────────────────────────────────────────────
% SLIDE 8 — USING THE NORMAL TABLE
% ─────────────────────────────────────────────────────────────────────────────
{
\setbeamercolor{background canvas}{bg=white}
\begin{frame}[t, plain]
  \begin{tikzpicture}[remember picture, overlay]
    \fill[navy] (current page.north west)
      rectangle ([yshift=-0.25cm]current page.north east);
  \end{tikzpicture}
  \vspace{0.45cm}\par
  \hspace{0.55cm}%
  \begin{minipage}{0.9\textwidth}

    \sectionlabel{Using the Normal Table (Table A)}

    \begin{columns}[T]
      \begin{column}{0.50\textwidth}
        \small
        \textbf{Table A gives:} $P(Z < z)$ = area to the \textbf{LEFT}.

        \vspace{0.5em}
        \begin{block}{Three cases}
          \small
          \renewcommand{\arraystretch}{1.4}
          \begin{tabular}{ll}
          $P(Z < z)$         & look up directly \\
          $P(Z > z)$         & $= 1 - P(Z < z)$ \\
          $P(z_1 < Z < z_2)$ & $= P(Z<z_2) - P(Z<z_1)$ \\
          \end{tabular}
        \end{block}

        \vspace{0.5em}
        \textbf{Calculator:}\\
        \texttt{normalcdf(lo, hi, $\mu$, $\sigma$)}\\
        \texttt{invNorm(area, $\mu$, $\sigma$)}
      \end{column}
      \begin{column}{0.46\textwidth}
        \small
        \textbf{Worked example — SAT $N(1060, 195)$:}\\
        Find $P(\text{SAT} < 900)$:
        \begin{enumerate}
          \setlength{\itemsep}{0.3em}
          \item $z = (900-1060)/195 \approx -0.82$
          \item Table A: $P(Z < -0.82) = 0.2061$
          \item $\approx 20.6\%$ scored below 900
        \end{enumerate}

        \vspace{0.4em}
        Find $P(\text{SAT} > 1400)$:
        \begin{enumerate}
          \setlength{\itemsep}{0.3em}
          \item $z = (1400-1060)/195 \approx 1.74$
          \item $P(Z > 1.74) = 1 - 0.9591 = 0.0409$
          \item $\approx 4.1\%$ scored above 1400
        \end{enumerate}
      \end{column}
    \end{columns}
  \end{minipage}
\end{frame}
}

% ─────────────────────────────────────────────────────────────────────────────
% SLIDE 9 — INVERSE NORMAL (PERCENTILE → VALUE)
% ─────────────────────────────────────────────────────────────────────────────
{
\setbeamercolor{background canvas}{bg=sky}
\begin{frame}[t, plain]
  \navyheader{Inverse Normal: Finding a Value from a Percentile}
  \vspace{0.8em}%
  \hspace{0.55cm}%
  \begin{minipage}{0.9\textwidth}
    \small
    \textbf{Question:} What SAT score is at the \textbf{90th percentile}?

    \vspace{0.4em}
    \textbf{Steps:}
    \begin{enumerate}
      \setlength{\itemsep}{0.5em}
      \item Find $z^*$ so that $P(Z < z^*) = 0.90$. \quad
            From Table A: $z^* \approx 1.28$
      \item Convert back: \quad $x = \mu + z^* \cdot \sigma = 1060 + 1.28(195) \approx \mathbf{1310}$
    \end{enumerate}

    \vspace{0.5em}
    A student scoring 1310 is at the 90th percentile — better than 90\% of test-takers.

    \vspace{0.5em}
    \begin{block}{General formula}
      \[
        x = \mu + z^* \cdot \sigma
      \]
      where $z^*$ = invNorm(percentile as a decimal).
    \end{block}

    \vspace{0.4em}
    \textbf{Try it:} What score is at the 25th percentile? ($z^* \approx -0.67$)
    \[
      x = 1060 + (-0.67)(195) \approx \mathbf{929}
    \]
  \end{minipage}
\end{frame}
}

% ─────────────────────────────────────────────────────────────────────────────
% SLIDE 10 — ASSESSING NORMALITY
% ─────────────────────────────────────────────────────────────────────────────
{
\setbeamercolor{background canvas}{bg=white}
\begin{frame}[t, plain]
  \begin{tikzpicture}[remember picture, overlay]
    \fill[navy] (current page.north west)
      rectangle ([yshift=-0.25cm]current page.north east);
  \end{tikzpicture}
  \vspace{0.45cm}\par
  \hspace{0.55cm}%
  \begin{minipage}{0.9\textwidth}

    \sectionlabel{Assessing Normality}

    \begin{columns}[T]
      \begin{column}{0.50\textwidth}
        \small
        \textbf{Method 1: Histogram}\\
        Look for a symmetric, unimodal, bell-shaped distribution
        with no strong skewness.

        \vspace{0.5em}
        \textbf{Method 2: Normal Probability Plot (NPP)}\\
        Plot each data value against its expected $z$-score if the
        data were Normal.
        \begin{itemize}
          \setlength{\itemsep}{0.3em}
          \item Points near a straight diagonal line
                $\Rightarrow$ approximately Normal
          \item Strong curves or bends $\Rightarrow$ skewness
          \item Individual outliers show as stragglers at the ends
        \end{itemize}

        \vspace{0.4em}
        \begin{block}{Key reminder}
          \small
          Always check before applying the Normal model.
          The model is an approximation --- use it only when reasonable.
        \end{block}
      \end{column}
      \begin{column}{0.46\textwidth}
        \begin{center}
        \textbf{Approximately Normal}\\[4pt]
        \begin{tikzpicture}[scale=0.72]
          \draw[thick] (0,0) -- (3.5, 3.5);
          \foreach \x/\y in {0.3/0.25, 0.8/0.85, 1.3/1.25, 1.8/1.82,
                              2.3/2.28, 2.8/2.75, 3.2/3.18} {
            \filldraw[navy] (\x, \y) circle (0.08);
          }
          \draw[thin, charcoal] (0,0) rectangle (3.5, 3.5);
          \node[below, font=\small\itshape] at (1.75, -0.2) {Points near line: Normal};
        \end{tikzpicture}

        \vspace{0.6em}
        \textbf{Not Normal (right skew)}\\[4pt]
        \begin{tikzpicture}[scale=0.72]
          \draw[thick] (0,0) -- (3.5, 3.5);
          \foreach \x/\y in {0.3/0.05, 0.8/0.30, 1.3/0.85, 1.8/1.80,
                              2.3/2.75, 2.8/3.20, 3.2/3.45} {
            \filldraw[navy] (\x, \y) circle (0.08);
          }
          \draw[thin, charcoal] (0,0) rectangle (3.5, 3.5);
          \node[below, font=\small\itshape] at (1.75, -0.2) {Curved: not Normal};
        \end{tikzpicture}
        \end{center}
      \end{column}
    \end{columns}
  \end{minipage}
\end{frame}
}

% ─────────────────────────────────────────────────────────────────────────────
% SLIDE 11 — GROUP ACTIVITY
% ─────────────────────────────────────────────────────────────────────────────
{
\setbeamercolor{background canvas}{bg=redbg}
\begin{frame}[t, plain]
  \navyheader{Group Activity: SAT Scores $N(1060, 195)$}
  \vspace{0.8em}%
  \hspace{0.55cm}%
  \begin{minipage}{0.9\textwidth}
    \small

    \textbf{Part 1 (8 min) --- Empirical Rule:}
    \begin{itemize}
      \setlength{\itemsep}{0.4em}
      \item Fill in the table: lower bound, upper bound, and \% for $\mu \pm 1,2,3\sigma$.
      \item What \% scored above 1255? Below 670? What score is the 84th percentile?
    \end{itemize}

    \vspace{0.5em}
    \textbf{Part 2 (8 min) --- $z$-Scores and Table A:}
    \begin{itemize}
      \setlength{\itemsep}{0.4em}
      \item Find $z$ for a score of 900 and find $P(\text{SAT} < 900)$.
      \item Find $P(\text{SAT} > 1400)$ and $P(900 < \text{SAT} < 1400)$.
    \end{itemize}

    \vspace{0.5em}
    \textbf{Part 3 (6 min) --- Inverse Normal and Unusual Values:}
    \begin{itemize}
      \setlength{\itemsep}{0.4em}
      \item Find the 90th percentile score.
      \item Is 1500 an unusual score? Justify with a $z$-score.
      \item Discuss: when is the Normal model \textit{not} appropriate?
    \end{itemize}
  \end{minipage}
\end{frame}
}

% ─────────────────────────────────────────────────────────────────────────────
% SLIDE 12 — EXIT TICKET
% ─────────────────────────────────────────────────────────────────────────────
{
\setbeamercolor{background canvas}{bg=sky}
\begin{frame}[t, plain]
  \navyheader{Exit Ticket}
  \vspace{1em}%
  \hspace{0.55cm}%
  \begin{minipage}{0.9\textwidth}
    \small
    \textit{Adult women's heights are approximately $N(65, 2.5)$ inches.}

    \vspace{0.5em}
    \begin{enumerate}
      \setlength{\itemsep}{0.7em}\setlength{\topsep}{0em}
      \item Find the $z$-score for a woman who is 70 inches tall.
            Interpret your answer in context.
      \item What proportion of women are taller than 70 inches?
            Show your work using Table A.
      \item Find the height at the 25th percentile.
            ($z^* \approx -0.67$.)
      \item Is 72.5 inches unusually tall? Justify with a $z$-score.
    \end{enumerate}

    \vspace{0.5em}
    \begin{block}{AP-Style}
      \small
      Newborn weights $\sim N(7.5, 1.1)$ lb. What proportion weigh between
      6.4 and 8.6 lb? Show work; interpret in context.
    \end{block}
  \end{minipage}
\end{frame}
}

% ─────────────────────────────────────────────────────────────────────────────
% SLIDE 13 — LESSON SUMMARY
% ─────────────────────────────────────────────────────────────────────────────
{
\setbeamercolor{background canvas}{bg=navy}
\begin{frame}[t, plain]
  \begin{tikzpicture}[remember picture, overlay]
    \fill[navylight] (current page.north west)
      rectangle ([yshift=-1.35cm]current page.north east);
  \end{tikzpicture}
  \vspace{1.5cm}
  \hspace{0.55cm}%
  \begin{minipage}{0.9\textwidth}
    {\color{goldacc}\bfseries\large Lesson 1.10 Summary}

    \vspace{0.5cm}
    {\color{white}\small
    \begin{itemize}
      \setlength{\itemsep}{0.5em}
      \item The \textbf{Normal distribution} $N(\mu, \sigma)$ is symmetric, unimodal,
            bell-shaped; fully described by mean and SD.
      \item The \textbf{68--95--99.7 Rule} gives approximate proportions within
            1, 2, 3 standard deviations of the mean.
      \item A \textbf{$z$-score} standardizes any value: $z = (x-\mu)/\sigma$.
            It counts standard deviations from the mean.
      \item \textbf{Table A} gives $P(Z < z)$ (area to the left).
            Subtract from 1 for ``greater than''; subtract areas for ``between''.
      \item Always \textbf{assess Normality} with a histogram or Normal probability plot
            before applying the Normal model.
    \end{itemize}
    }

    \vspace{0.6cm}
    {\color{skymid}\small\itshape
      Next: Unit 2 --- Exploring Two-Variable Data (scatterplots, correlation, regression).}
  \end{minipage}
\end{frame}
}

\end{document}
